\begin{table*}[t]
\centering
\small
\renewcommand{\arraystretch}{1.1}
\scalebox{0.85}{
\begin{tabular}{lcccccc|cccccc}
\toprule
\multicolumn{1}{c}{} &
\multicolumn{6}{c}{\textbf{Short-term} ($50^{th}\text{frame}$)} &
\multicolumn{6}{c}{\textbf{Long-term} ($\geq 200^{th}\text{frame}$)}\\
\cmidrule(lr){2-7}\cmidrule(lr){8-13}
\multicolumn{1}{c}{} &
\textbf{LPIPS} $\downarrow$ &
\textbf{PSNR} $\uparrow$ &
\textbf{SSIM} $\uparrow$ &
\textbf{FID} $\downarrow$ &
$R_{\text{dist}}$ $\downarrow$ &
$T_{\text{dist}}$ $\downarrow$ &
\textbf{LPIPS} $\downarrow$ &
\textbf{PSNR} $\uparrow$ &
\textbf{SSIM} $\uparrow$ &
\textbf{FID} $\downarrow$ &
$R_{\text{dist}}$ $\downarrow$ &
$T_{\text{dist}}$ $\downarrow$ \\
\midrule
   GeoGPT~\cite{geogpt}
    & 0.444 & 13.35  & --- & 26.72 & ---&---
    & 0.674 & 9.54  & --- & 41.87 & ---& ---\\
   Lookout~\cite{ren2022look}
    & 0.378 & 14.43  & --- &28.86  & 1.241& 0.876
    & 0.658 & 10.51  & --- &  58.12& 1.142&0.924\\
   PhotoNVS~\cite{photometricnvs}
    &  0.333& 15.51  & ---&21.76  & ---& ---
    & 0.588 &  11.54 & --- & 41.95 &--- &--- \\


   GenWarp~\cite{seo2024genwarp}
    & 0.436 &12.03   & 0.144& 29.69 & 0.564& 0.059 
    & 0.613 & 9.56  & 0.085 &  36.40& 0.850& \textbf{0.251} \\
   MotionCtrl~\cite{wang2024motionctrl}
    & 0.424  & 12.00   & 0.148 & 19.25 & 0.341 & 0.356
    & 0.605 & 9.13  & 0.083 & 35.45 & \textbf{0.748} & 0.609\\

   ViewCrafter~\cite{yu2024viewcrafter}
    & 0.377  & 16.97  & 0.262 & 25.39 &1.562 &0.208
    &  0.592 & 9.74  & 0.148 &  34.12& 2.177& 0.827\\


    SEVA~\cite{seva} ($K=17$)
    &  0.293 & 18.33  & 0.382  & 17.29 & 0.223& 0.118
    & 0.455 & 13.98  & 0.216 & 23.82 & 1.125 & 0.742\\

    \midrule

    \mname~($K=4$)
    & \textbf{0.287} & \textbf{18.49}  & \textbf{0.406} &\textbf{17.12}  & \textbf{0.219}& \textbf{0.039}
    & 0.493 & 13.12 & 0.183 &  27.15& 0.811& 0.499 \\

    \mname~($K=17$)
    & 0.293& 18.33  & 0.382 & 17.29 &0.223 & 0.118
    & \textbf{0.452}& \textbf{14.09} &  \textbf{0.227} & \textbf{23.56}  & 0.982 & 0.432 \\
\bottomrule
\end{tabular}
}
\caption{\textbf{Single-view NVS on RealEstate10K.}
Frames are subsampled with 10-frame intervals.
$K$ denotes the number of context views.
In the short-term setting, all past frames fit in the context windows for both SEVA~\cite{seva} and \mname; hence, SEVA~\cite{seva} and \mname~($K=17$) yield identical results since \mname~uses SEVA as the backbone.
\mname~($K=4$) is fine-tuned on RealEstate10K for this setting, yielding better performance than SEVA~\cite{seva} and \mname~($K=17$) in short-term scenarios.
\mname~surpasses baselines on most metrics in this benchmark. However, its \textbf{true advantage—spatial consistency—is not fully reflected} because RealEstate10K trajectories rarely revisit previously observed areas.
We address this limitation in~\cref{exp:cycle}.
}%
\label{tab:long-term}
\vspace{-1mm}
\end{table*}

\begin{table}[t]
\centering
\setlength{\tabcolsep}{3pt}
\scalebox{0.76}{
\begin{tabular}{lcccccc}
\toprule
\textbf{Method} & \textbf{LPIPS}~$\downarrow$ & \textbf{PSNR}~$\uparrow$  & \textbf{SSIM}~$\uparrow$& \textbf{FID}~$\downarrow$  & $R_{\text{dist}}$ $\downarrow$ & $T_{\text{dist}}$ $\downarrow$ \\
\midrule
%
%
Look-out~\cite{ren2022look}
&  0.809&   8.41&  0.069&  38.34&  1.262&0.341\\
%
%
GenWarp~\cite{seo2024genwarp}
&  0.507&   11.13&  0.134&  32.94&  0.848&0.241\\

MotionCtrl~\cite{wang2024motionctrl}
& 0.589 & 9.07  &  0.096& 26.86 &  0.871&0.293\\

ViewCrafter~\cite{yu2024viewcrafter}
&  0.401&   11.82&  0.217& 24.72 & 0.902&  0.492\\

SEVA~\cite{seva} ($K=17$)
&  0.401&   11.82&  0.217& 24.72 & 0.902&  0.492\\

\midrule

\mname~($K=4$)&  0.397&  15.72&  0.297& 24.97&  \textbf{0.821} & 0.392\\

\mname~($K=17$)&  \textbf{0.304}&  \textbf{18.15}&  \textbf{0.377}& \textbf{24.18}&  0.892& \textbf{0.165}\\

\bottomrule
\end{tabular}
}
\caption{\textbf{Results on cycle trajectories from RealEstate10K.} 
%
%
}%
\label{tab:cycle}
\end{table}

\begin{table}[t]
    \centering
    \setlength{\tabcolsep}{5pt}
    \scalebox{0.76}{
    \begin{tabular}{lccccc}
    \toprule
    \textbf{Method} & \textbf{LPIPS}~$\downarrow$ & \textbf{PSNR}~$\uparrow$ & \textbf{SSIM}~$\uparrow$  & $R_{\text{dist}}$ $\downarrow$ & $T_{\text{dist}}$ $\downarrow$ \\
    \midrule
   %
   %
   Look-out~\cite{ren2022look}
    &  0.792&   8.52&  0.008&  \textbf{0.727}&  1.499\\
   %
   %
   GenWarp~\cite{seo2024genwarp}
    &  0.521&   10.27&  \textbf{0.129}& 0.785&  1.982\\

   MotionCtrl~\cite{wang2024motionctrl}
    &  0.692&   8.21&  0.082& 0.842& \textbf{0.129}\\

   ViewCrafter~\cite{yu2024viewcrafter}
       &0.494& 13.62 & \textbf{0.129}  &1.643 & 0.492 \\

    \midrule

   \mname~($K=4$)
    & \textbf{0.472} & \textbf{14.11} & 0.121 & 1.204 &0.387 \\

    \bottomrule
    \end{tabular}
    }
    \caption{\textbf{Results on cycle trajectories from Tanks and Temples.} 
    %
    }%
    \label{tab:tank}
    \vspace{0mm}
\end{table}




\section{Experiments}%
\label{sec:experiments}

We evaluate our method on established benchmarks for camera-conditioned autoregressive view generation from a single image, comparing our models to previous open-source approaches (\cref{exp:short-term}).
These benchmarks, however, have a critical limitation: camera trajectories rarely revisit previously observed regions, unlike real-world behavior where humans or robots often return to or look at previously explored areas.
To address this limitation, we propose a cycle-trajectory evaluation protocol that extends these benchmarks by making the camera return to the starting position along the same path in reverse order (\cref{exp:cycle}).
Additionally, we collect a small dataset with trajectories that revisit observed regions for qualitative evaluation in~\cref{fig:more_demo}.
Finally, we conduct an ablation study in \cref{exp:ablation}.

\subsection{Experimental setup}%
\label{sec:experimental_setup}

\paragraph{Implementation details.}%
We use pre-trained SEVA~\cite{seva} as the generator $\psi$.
SEVA uses a fixed total number of reference and target frames ($K+M=21$).
For our experiments, we use $M=4$ target views, which leaves $K=17$ context views.
However, this configuration demands substantial computational resources that may be prohibitive for many applications.
To address this limitation, we fine-tune a more efficient version with LoRA~\cite{hu22lora:} on the RealEstate10K training split~\cite{realestate} that works with a reduced number of reference frames ($K=4$) and target views ($M=4$), and use this efficient version for our main experiments.
We demonstrate (\cref{exp:cycle,exp:ablation}) that this fine-tuned version, when combined with~\mname, achieves comparable performance to the original SEVA while delivering approximately $12\times$ speedup at inference time.
We provide more technical details in~\acref{appen:implementation}.

\paragraph{Datasets.}%
\label{sec:datasets}
We consider two real-world datasets for novel view generation evaluation: RealEstate10K~\cite{realestate}, comprising indoor scene video clips, and Tanks-and-Temples~\cite{knapitsch2017tanks}, including indoor and outdoor scenes with larger camera motions.
For qualitative evaluation, we use in-the-wild images collected from the internet or captured with phone cameras as input views (see \cref{fig:more_demo}).

\paragraph{Evaluation metrics.}%
\label{sec:metrics}
We evaluate generated novel view quality based on three factors:
(1) overall quality of generated novel views, comparing generated novel view distributions to test set distributions using Fréchet Image Distance (FID)~\cite{fid};
(2) the ability of the model to preserve image details across views, measuring peak signal-to-noise ratio (PSNR) for pixel differences, along with LPIPS~\cite{lpips} and SSIM~\cite{ssim}, as in~\cite{ren2022look}; and
(3) alignment between generated camera poses and ground truth, following~\cite{he2024cameractrl}, where we express estimated camera poses relative to the first frame and normalize translation by the furthest frame.
We use DUSt3R~\cite{dust3r_cvpr24} to extract poses from generated views.
We calculate rotation distance $R_{\text{dist}}$ by comparing ground truth and extracted rotation matrices of each generated sequence:
\[
R_{\text{dist}} = \arccos\left(0.5\left(\text{tr}(\mathbf{R}_{\text{gen}} \mathbf{R}^{T}_{\text{gt}}) - 1\right)\right),
\]
where $\mathbf{R}$ denotes a rotation matrix.
We compute translation distance $t_{\text{dist}}$ as:
\[
t_{\text{dist}} = \left \| \mathbf{t}_{\text{gt}} - \mathbf{t}_{\text{gen}} \right \|_2.
\]

\begin{figure*}[ht]
\centering
\includegraphics[width=0.99\textwidth]{imgs/vmem_long-term_comparison_figure.pdf}
%
\caption{\textbf{Qualitative comparison of the cycle trajectory frames} ($\geq 400$ frames). \mname~autoregressively generates new views while maintaining memory of all previously generated frames, ensuring consistency when revisiting previously seen locations.}%
\label{fig:longqualitative_comparison}
\end{figure*}


\subsection{Short- and long-term view generation}%
\label{exp:short-term}

We first evaluate \mname~on the benchmark from~\cite{ren2022look, photometricnvs}, which defines two scenarios: \emph{short-term} and \emph{long-term} novel view generation.
Beginning from the first frame of each ground-truth test sequence, the model autoregressively generates images along the ground truth camera trajectories with 10-frame intervals.
The \emph{short-term} scenario evaluates the fifth generated image (50 frames from the initial view), while the \emph{long-term} scenario evaluates the final image ($\geq$200 frames from the initial view).
Following the protocol, we only consider samples with more than 200 frames for the \emph{long-term} scenario and more than 50 frames for the \emph{short-term} scenario.
We report results in \cref{tab:long-term}.
\mname~substantially outperforms all baselines on the key metrics, demonstrating superior scene extrapolation.
However, the benchmark trajectories rarely revisit observed regions, limiting evaluation of long-term scene consistency.

\begin{table}[]
\centering
\setlength{\tabcolsep}{2.5pt}
\scalebox{0.682}{
\begin{tabular}{clccccc}
\toprule
\makecell{\textbf{Context} \\ \textbf{views $K$}}& \textbf{View Retrieval} & \textbf{LPIPS}~$\downarrow$ & \textbf{PSNR}~$\uparrow$ & \textbf{SSIM}~$\uparrow$ & $R_{\text{dist}}$~$\downarrow$ & $T_{\text{dist}}$~$\downarrow$ \\
\midrule
\multirow{4}{*}{\textbf{17}} 
& Temporal         & 0.477 & 13.92 & 0.188 & 0.976 & 0.254 \\
& Camera Distance (SEVA~\cite{seva})  & 0.397 & 15.72 & 0.297 & \textbf{0.821} & 0.392 \\
& Field of View & 0.374 & 15.75 & 0.292 & 0.911 & 0.382 \\
& \mname             & \textbf{0.304} & \textbf{18.15} & \textbf{0.377} & 0.892 & \textbf{0.165} \\
\midrule
\multirow{4}{*}{\textbf{4}} 
& Temporal         & 0.794 & 7.52  & 0.018 & 1.942 & 0.458 \\
& Camera Distance  & 0.422 & 13.27 & 0.187 & 1.787 & 0.319 \\
& Field of View & 0.424 & 13.11 & 0.192 & 1.782 & 0.285 \\
& \mname             & \textbf{0.381} & \textbf{14.82} & \textbf{0.275} & \textbf{0.793} & \textbf{0.124} \\
\bottomrule
\end{tabular}
}
\caption{\textbf{Ablation study on cycle trajectories from RealEstate10K.} We compare different view retrieval strategies for two context view counts $K$
%
}%
\label{tab:ablation}
\vspace{-4mm}
\end{table}


\subsection{Long-term view generation with revisitations}%
\label{exp:cycle}

To evaluate long-term scene consistency when revisiting observed regions—a scenario that rarely occurs in existing benchmarks—we propose a cycle-trajectory evaluation protocol as a proxy for real revisitation scenarios in human and robotic exploration.
In this protocol, the model generates frames following trajectories from initial to final poses, then returns along the same path to the starting position.
We evaluate metrics every 10 frames on the return trajectories, using all test sequences.
Following this protocol, we report quantitative results in \cref{tab:cycle} and qualitative comparisons in \cref{fig:longqualitative_comparison}.
\mname~outperforms all baselines across all metrics, demonstrating superior visual consistency when generating previously observed regions.

To validate generalization, we evaluate on Tanks-and-Temples containing indoor and outdoor scenes with more dynamic camera trajectories.
We use all six advanced scenes and apply the same \emph{cycle-trajectory} protocol using the first 50 frames, evaluating every frame without temporal subsampling.
Results are in \cref{tab:tank}.
\mname~consistently outperforms all baselines across in-domain (\cref{tab:cycle}) and out-of-domain (\cref{tab:tank}) evaluations on the majority of metrics.

We also qualitatively evaluate \mname~across diverse in-the-wild scenes in \cref{fig:more_demo}, demonstrating generation of diverse, consistent, high-quality long videos.

\subsection{Ablation study}%
\label{exp:ablation}


We ablate our surfel-based memory indexing for context view retrieval in \mname.
We evaluate our approach in the cycle-trajectory setting against three alternative retrieval strategies:
(1) temporal-based retrieval, selecting the most recent $K$ views from memory,
(2) camera distance-based retrieval, selecting top $K$ views with cameras closest to the target view camera—also used by SEVA~\cite{seva}, our backbone generator, and
(3) field-of-view-based retrieval, selecting top $K$ views with highest field-of-view overlap with the target view.
We also evaluate performance trade-offs between our lightweight, fine-tuned generator using $K=4$ context views versus original SEVA with $K=17$ context views.
We report results in \cref{tab:ablation}.

Our findings show \mname~consistently improves generator performance across both numbers of context views.
Improvement is most pronounced with fewer context views ($K=4$), highlighting the effectiveness of our surfel-based memory index.

\paragraph{Computational efficiency.}
A key practical advantage of our approach is its computational efficiency.
\mname~with only $K=4$ views achieves $12\times$ speed improvement over the original SEVA while recovering most of the performance achieved with $K=17$ views.
On an RTX 4090 GPU, our method generates frames in 4.2 seconds compared to 50 seconds for SEVA, representing a step towards real-time interactive scene exploration.
